%% kapa: exam=1 round=33 year=2022 period=1 form=A subject=civil_law subject_ko=민법 seq=10 number=10
%% source: Q-Net 감정평가사 기출문제 (cst00309) artlSeq=5209997 file=raw/1008/5209997/2022년도 제33회 감정평가사 1차시험 1교시(민법,경제학,부동산학).pdf
%% conversion: kordoc 4.24.51 -> tools/gen1.py (auto); review=proofread-sonnet; flags=none
\begin{question}{10}
물건에 관한 설명으로 옳지 않은 것은? (다툼이 있으면 판례에 따름)
\begin{choices}
  \item 주물에 대한 압류의 효력은 특별한 사정이 없는 한 종물에는 미치지 않는다.
  \item 사람의 유골은 매장·관리의 대상이 될 수 있는 유체물이다.
  \item 전기 기타 관리할 수 있는 자연력은 물건이다.
  \item 법정과실은 수취할 권리의 존속기간 일수의 비율로 취득함이 원칙이다.
  \item 주물만 처분하고 종물은 처분하지 않기로 하는 특약은 유효하다.
\end{choices}
\end{question}
\answer{1}
